\section{GPT-3、The Pile、Gopher、LLaMA、RefinedWeb}

GPT-3 数据包括 processed Common Crawl、WebText2、神秘 books corpora 和 Wikipedia，总计约 400B tokens。处理上使用 quality classifier 区分 WebText/Wikipedia/Books 与其他网页，并做 fuzzy deduplication。

The Pile 是开源模型社区对 GPT-3 数据不透明的回应，聚合 22 个高质量 domain，包括 Pile-CC、PubMed Central、arXiv、Enron emails、Project Gutenberg、Books3、StackExchange 等。

\begin{figure}[H]
\centering
\includegraphics[width=0.76\textwidth]{the-pile.png}
\caption{The Pile：22 个 curated high-quality domains 的组合。}
\figsource{Stanford CS324 lecture image，本地化后供 CS336 Lecture 13 使用。}
\end{figure}

\begin{knowledgebox}{读图：The Pile 的价值和争议}
The Pile 让数据组成更透明，推动开源模型训练。但其中 Books3 来自 shadow library Bibliotik，后来因 copyright lawsuits 下线。透明不等于没有法律/伦理风险。
\end{knowledgebox}

Gopher 的 MassiveText 说明数据描述本身也很有价值：MassiveWeb、C4、Books、News、GitHub、Wikipedia，并使用 English filtering、deduplication、train-test overlap、manual quality rules 和 Google SafeSearch toxicity filtering。LLaMA 则使用 CommonCrawl/CCNet/C4/GitHub/Wikipedia/Project Gutenberg/Books3/arXiv/StackExchange，总计约 1.2T tokens，并被 RedPajama/SlimPajama 复刻。

RefinedWeb 的主张是 web data is all you need：使用 WARC 和 trafilatura 抽取文本，Gopher rules 过滤，避免 ML-based filtering 以减少 bias，并用 MinHash over 5-grams 做 fuzzy deduplication。FineWeb 在此基础上扩展到 95 个 Common Crawl dumps，结果达到 15T tokens。

\begin{knowledgebox}{这些数据集的代际变化}
早期数据集强调“有一批可训练文本”：Wikipedia、BooksCorpus、WebText。GPT-3 之后，数据集开始强调“大规模 web + 多源混合 + 去重过滤”：The Pile、MassiveText、LLaMA data、RefinedWeb。到 Dolma/DCLM/Nemotron-CC 阶段，重点进一步变成“过滤器和数据处理算法本身”：谁能从巨大 raw pool 中筛出更高质量、更少偏差、更可审计的数据。
\end{knowledgebox}

\begin{warningbox}{不要把 token 数当成唯一质量指标}
15T tokens 不必然优于 3T tokens，3T tokens 也不必然优于 800GB。数据价值取决于覆盖、质量、去重、许可证、语言分布、代码/数学/学术比例、污染控制和与目标能力的匹配。Scaling laws 只能在数据分布相对稳定时外推；换数据 pipeline 会改变曲线。
\end{warningbox}

\begin{warningbox}{手工规则和模型过滤的偏差不同}
手工规则透明、便宜、可解释，但可能粗糙，例如删除包含代码符号的页面。模型过滤更强，但会把训练 classifier 的正负例偏差扩散到整个语料。数据过滤没有中立方案，只有可审计的 trade-off。
\end{warningbox}

